\documentclass{article}
\usepackage[margin=1in]{geometry}
\usepackage{hyperref}
\usepackage{enumitem}

\setlist[itemize]{leftmargin=*, itemsep=2pt, topsep=2pt}

\title{Unified Content Template: Experience \& Projects}
\author{Template for Human/AI Content Drafting}
\date{Last Updated: 2026-02-18}

\begin{document}

\maketitle

\section{Purpose}
This single file is the canonical template for drafting detailed source content entries used to generate tailored resumes and cover letters.

Use this template to write:\
(1) work experience entries and\
(2) project entries.\
The output should be truthful, specific, and quantifiable.

\section{How to Use This Template}
\begin{itemize}
    \item Copy this file, rename it to match the target entry, and fill in placeholders in brackets.
    \item Keep claims accurate and consistent with known background materials.
    \item Prefer measurable impact (time saved, metric lift, error reduction, cost reduction, ranking, etc.).
    \item For ML entries, use investigation-first narrative where possible: \textbf{Problem complexity $\rightarrow$ Investigation/Hypothesis $\rightarrow$ Experiment/Decision $\rightarrow$ Impact}.
    \item Mention constraints explicitly (data scarcity, class imbalance, domain shift, latency/cost limits).
    \item Hardware wording: use \textbf{A100 GPUs} by default; mention TPU only for Kaggle competition projects where TPU was actually used.
\end{itemize}

\section{Expected Input Format (for AI filling this template)}
Provide structured input before drafting content. Recommended fields:

\subsection*{Global Metadata}
\begin{itemize}
    \item \textbf{Name of entry:} [Required]
    \item \textbf{Entry type:} [experience | project]
    \item \textbf{One-line summary:} [Required]
    \item \textbf{Priority for resume use:} [high | medium | low]
\end{itemize}

\subsection*{Evidence and Constraints}
\begin{itemize}
    \item \textbf{Known facts:} [bullet list of verifiable facts]
    \item \textbf{Unknowns/placeholders:} [items requiring follow-up]
    \item \textbf{Numbers and metrics available:} [list all quantitative evidence]
    \item \textbf{No-claim zones:} [areas where details are uncertain and must not be overstated]
\end{itemize}

\section{Experience Entry Template}

\subsection{Header}
\begin{itemize}
    \item \textbf{Company:} [Company Name]
    \item \textbf{Title:} [Role Title]
    \item \textbf{Dates:} [Month Year -- Month Year]
    \item \textbf{Location:} [City/Remote]
    \item \textbf{Team/Org:} [Team details, if known]
    \item \textbf{Domain:} [e.g., CV inspection, NLP, BI]
\end{itemize}

\subsection{Company/Org Context}
[2--4 sentences: What company does, what team mission was, why this work mattered.]

\subsection{Role Scope and Ownership}
\begin{itemize}
    \item [What you owned end-to-end]
    \item [Who you collaborated with]
    \item [Decision scope: model/data/evaluation/deployment]
\end{itemize}

\subsection{Key Initiative 1: [Initiative Name]}
\subsubsection*{Problem Statement}
[Describe baseline state, pain point, target users, and success criteria.]

\subsubsection*{Investigation and Hypothesis}
\begin{itemize}
    \item [Failure analysis performed]
    \item [Alternative approaches considered]
    \item [Why selected direction was chosen]
\end{itemize}

\subsubsection*{Technical Design}
\begin{itemize}
    \item \textbf{Task framing:} [classification/detection/ranking/etc.]
    \item \textbf{Modeling approach:} [models, training strategy]
    \item \textbf{Data strategy:} [labeling, sampling, augmentation, splits]
    \item \textbf{Evaluation design:} [primary metrics and why]
    \item \textbf{Infra/tools:} [PyTorch/TensorFlow, tracking, compute]
\end{itemize}

\subsubsection*{Results and Impact}
\begin{itemize}
    \item \textbf{Performance outcome:} [metric before $\rightarrow$ after]
    \item \textbf{Business/operational impact:} [time/cost/quality impact]
    \item \textbf{Deployment status:} [prototype | production | sunset + reason]
\end{itemize}

\subsubsection*{Challenges and Tradeoffs}
\begin{itemize}
    \item [Constraint + mitigation]
    \item [Tradeoff accepted and rationale]
\end{itemize}

\subsection{Key Initiative 2: [Initiative Name] (Optional)}
[Repeat the same structure as Initiative 1.]

\subsection{Additional Contributions (Optional)}
\begin{itemize}
    \item [Mentorship, hiring, code quality, process improvements]
    \item [Cross-functional influence or internal standards]
\end{itemize}

\subsection{Experience Summary for Resume Selection}
\begin{itemize}
    \item \textbf{Best bullets to extract:} [3--6 concise bullets]
    \item \textbf{Best-fit role types:} [MLE, applied scientist, research engineer, etc.]
    \item \textbf{When to de-prioritize:} [conditions where this entry is less relevant]
\end{itemize}

\section{Project Entry Template}

\subsection{Header}
\begin{itemize}
    \item \textbf{Project name:} [Project/Competition Name]
    \item \textbf{Type:} [Kaggle | research | side project | production prototype]
    \item \textbf{Timeline:} [Month Year -- Month Year]
    \item \textbf{Team setup:} [individual/team size]
    \item \textbf{Link(s):} [repo, paper, competition URL]
\end{itemize}

\subsection{Problem and Scope}
\begin{itemize}
    \item \textbf{Objective:} [What the system must do]
    \item \textbf{Scale:} [data size, classes, throughput, users]
    \item \textbf{Main constraints:} [compute/data/latency/fairness/etc.]
\end{itemize}

\subsection{Approach}
\begin{itemize}
    \item \textbf{Baselines:} [what was tried first]
    \item \textbf{Core method:} [final architecture + training approach]
    \item \textbf{Data/evaluation strategy:} [splits, metrics, ablations]
    \item \textbf{System implementation:} [pipeline, serving, optimization]
\end{itemize}

\subsection{Results}
\begin{itemize}
    \item \textbf{Primary outcomes:} [rank/metric/throughput gains]
    \item \textbf{Comparative improvement:} [baseline $\rightarrow$ final]
    \item \textbf{What did not work:} [failed ideas and lessons]
\end{itemize}

\subsection{Technical Stack}
\begin{itemize}
    \item \textbf{Languages:} [Python/SQL/etc.]
    \item \textbf{Frameworks:} [PyTorch/TensorFlow/etc.]
    \item \textbf{Infra:} [A100/cloud/local; TPU only if true for this project]
    \item \textbf{Tooling:} [tracking/versioning/experiment management]
\end{itemize}

\subsection{Project Summary for Resume Selection}
\begin{itemize}
    \item \textbf{Top resume bullets:} [2--4 bullets]
    \item \textbf{Role alignment:} [which job types this project supports]
    \item \textbf{Key keyword coverage:} [CV/NLP/OOD/fairness/scaling/etc.]
\end{itemize}

\section{Writing Quality Checklist}
Before finalizing, verify:
\begin{itemize}
    \item Every major claim has supporting evidence.
    \item Numbers are concrete and internally consistent.
    \item No over-claiming; uncertainty is marked explicitly.
    \item Problem-method-impact flow is clear.
    \item Content is reusable for selective resume extraction.
\end{itemize}

\section{Minimal Example Snippet (Format Reference)}
\subsection*{Experience Initiative Example (short)}
\begin{itemize}
    \item \textbf{Problem:} Manual review volume was 500K images/season; baseline detector produced 60\% FPR@95.
    \item \textbf{Investigation:} Error analysis showed localization quality was not the business bottleneck; false alarms were.
    \item \textbf{Decision:} Reframed task to image-level classification, optimized for FPR@95, used CLIP fine-tuning with class-imbalance augmentation.
    \item \textbf{Impact:} Reduced FPR@95 to 20\%, enabling 80\% auto-rejection of negatives and major review-time savings.
\end{itemize}

\subsection*{Project Result Example (short)}
\begin{itemize}
    \item \textbf{Outcome:} Ranked 30/542 in large-scale image retrieval challenge.
    \item \textbf{System note:} Built scalable training pipeline for 100GB+ data with efficient mixed-precision training and retrieval indexing.
\end{itemize}

\end{document}
